\documentclass[10pt]{article}
\usepackage[utf8]{inputenc}
\usepackage[vietnamese]{babel}
\usepackage{amsmath,amssymb,amsfonts}
\usepackage{geometry}
\geometry{
    paperwidth=8.5in,
    paperheight=11in,
    top=1.35cm,
    bottom=1.8cm,
    left=1.35cm,
    right=1.35cm,
    headheight=14pt,
    footskip=25pt
}
\usepackage{xcolor}
\usepackage{colortbl}
\usepackage{tikz}
\usepackage{fancyhdr}
\usepackage{calc}
\usepackage{enumitem}

% Bảng màu chuẩn giáo trình Stewart Calculus
\definecolor{stewartcyan}{RGB}{0, 163, 218}
\definecolor{stewarttableblue}{RGB}{227, 238, 247}
\definecolor{stewartfooter}{RGB}{100, 100, 100}

\pagestyle{fancy}
\fancyhf{}
\renewcommand{\headrulewidth}{0pt}
\fancyhead[L]{%
    \fontsize{9}{11}\selectfont%
    \textbf{\textsf{342}}\hspace{22pt}%
    \textbf{\textsf{CHƯƠNG 4}}\hspace{10pt}%
    \textsf{Các ứng dụng của đạo hàm}%
}

\fancyfoot[C]{%
    \fontsize{6.5}{8}\selectfont\color{stewartfooter}%
    Bản quyền thuộc Cengage Learning. Mọi quyền được bảo lưu. Không được sao chép, quét hoặc nhân bản toàn bộ hay một phần. WCN 02-200-203\\
    Do quyền điện tử, một số nội dung của bên thứ ba có thể bị ẩn khỏi sách điện tử và/hoặc chương điện tử. Đánh giá biên tập cho thấy rằng mọi nội dung bị lược bỏ không ảnh hưởng đáng kể đến trải nghiệm học tập tổng thể. Cengage Learning có quyền gỡ bỏ nội dung bổ sung vào bất kỳ lúc nào nếu các hạn chế bản quyền sau đó yêu cầu.%
}

\setlength{\parindent}{0pt}
\setlength{\parskip}{0pt}

\begin{document}

% ==============================================================================
% PHẦN ĐẦU TRANG: LỜI GIẢI TIẾP NỐI VÍ DỤ TỪ TRANG TRƯỚC (CỘT NỘI DUNG CHÍNH)
% ==============================================================================
\noindent
\hspace*{155pt}%
\begin{minipage}{380pt}
    \fontsize{9.5}{13.5}\selectfont
    với mỗi đơn vị bán thêm, mức giảm giá sẽ là $\frac{1}{20} \times 10$ và hàm cầu là
    \[
    p(x) = 350 - \frac{10}{20}(x - 200) = 450 - \frac{1}{2}x
    \]
    Hàm doanh thu là
    \[
    R(x) = xp(x) = 450x - \frac{1}{2}x^2
    \]
    Vì $R'(x) = 450 - x$, ta thấy rằng $R'(x) = 0$ khi $x = 450$. Giá trị này của $x$ mang lại giá trị cực đại tuyệt đối theo Kiểm tra đạo hàm cấp một (hoặc đơn giản bằng cách nhận thấy đồ thị của $R$ là một parabol có bề lõm quay xuống). Mức giá tương ứng là
    \[
    p(450) = 450 - \frac{1}{2}(450) = 225
    \]
    và mức giảm giá khuyến mãi là $350 - 225 = 125$. Do đó, để tối đa hóa doanh thu, cửa hàng nên đưa ra mức giảm giá là \$125.\hfill{\color{stewartcyan}\rule{1.1ex}{1.1ex}}
\end{minipage}

\vspace{20pt}

% ==============================================================================
% THANH TIÊU ĐỀ: 4.7 BÀI TẬP (EXERCISES)
% ==============================================================================
\noindent
\begin{tikzpicture}[x=1pt,y=1pt]
    \node[inner sep=0pt, outer sep=0pt, anchor=west] (heading) at (46, 0) {%
        {\fontsize{16}{19}\selectfont\sffamily\bfseries\color{stewartcyan}4.7}\hspace{8pt}%
        {\fontsize{16}{19}\selectfont\sffamily\bfseries Bài tập}%
    };
    
    \pgfmathsetmacro{\ymid}{5.5}
    \pgfmathsetmacro{\ygap}{2.3}
    
    % 3 đường kẻ cyan bên trái
    \draw[stewartcyan, line width=0.6pt] (0, \ymid + \ygap) -- (38, \ymid + \ygap);
    \draw[stewartcyan, line width=0.6pt] (0, \ymid) -- (38, \ymid);
    \draw[stewartcyan, line width=0.6pt] (0, \ymid - \ygap) -- (38, \ymid - \ygap);
    
    % 3 đường kẻ cyan bên phải
    \path (heading.east);
    \pgfgetlastxy{\hx}{\hy};
    \draw[stewartcyan, line width=0.6pt] (\hx + 8pt, \ymid + \ygap) -- (\textwidth, \ymid + \ygap);
    \draw[stewartcyan, line width=0.6pt] (\hx + 8pt, \ymid) -- (\textwidth, \ymid);
    \draw[stewartcyan, line width=0.6pt] (\hx + 8pt, \ymid - \ygap) -- (\textwidth, \ymid - \ygap);
\end{tikzpicture}

\vspace{12pt}

% ==============================================================================
% BỐ CỤC 2 CỘT CHO BÀI TẬP
% ==============================================================================
\noindent
% -------------------------- CỘT TRÁI --------------------------
\begin{minipage}[t]{0.48\textwidth}
\fontsize{9}{12.2}\selectfont

\hangindent=13pt\hangafter=1
\textbf{1.} Xét bài toán sau: tìm hai số có tổng bằng 23 và tích của chúng đạt giá trị lớn nhất.
\begin{enumerate}[label=\textbf{(\alph*)}, leftmargin=18pt, labelsep=4pt, itemsep=2pt, parsep=0pt, topsep=2.5pt]
    \item Lập một bảng giá trị, tương tự như bảng bên dưới, sao cho tổng của các số trong hai cột đầu tiên luôn bằng 23. Trên cơ sở các số liệu trong bảng của bạn, hãy ước lượng đáp số cho bài toán.
    \item Sử dụng phép tính vi tích phân để giải bài toán và so sánh với câu trả lời của bạn ở phần (a).
\end{enumerate}

\vspace{4pt}
\begin{center}
\renewcommand{\arraystretch}{1.15}
\setlength{\tabcolsep}{9pt}
{\fontsize{8.5}{11}\selectfont
\begin{tabular}{|c|c|c|}
    \hline
    \rowcolor{stewarttableblue}
    \textsf{Số thứ nhất} & \textsf{Số thứ hai} & \textsf{Tích} \\
    \hline
    1 & 22 & 22 \\
    2 & 21 & 42 \\
    3 & 20 & 60 \\
    \vdots & \vdots & \vdots \\
    \hline
\end{tabular}}
\end{center}
\vspace{5pt}

\hangindent=13pt\hangafter=1
\textbf{2.} Tìm hai số có hiệu bằng 100 và tích của chúng đạt giá trị nhỏ nhất.

\vspace{6pt}
\hangindent=13pt\hangafter=1
\textbf{3.} Tìm hai số dương có tích bằng 100 và tổng của chúng đạt giá trị nhỏ nhất.

\vspace{6pt}
\hangindent=13pt\hangafter=1
\textbf{4.} Tổng của hai số dương là 16. Giá trị nhỏ nhất có thể có của tổng các bình phương của chúng là bao nhiêu?

\vspace{6pt}
\hangindent=13pt\hangafter=1
\textbf{5.} Khoảng cách thẳng đứng lớn nhất giữa đường thẳng $y = x + 2$ và parabol $y = x^2$ với $-1 \leqslant x \leqslant 2$ là bao nhiêu?

\vspace{6pt}
\hangindent=13pt\hangafter=1
\textbf{6.} Khoảng cách thẳng đứng nhỏ nhất giữa các parabol $y = x^2 + 1$ và $y = x - x^2$ là bao nhiêu?

\vspace{6pt}
\hangindent=13pt\hangafter=1
\textbf{7.} Tìm kích thước của một hình chữ nhật có chu vi $100\text{ m}$ sao cho diện tích của nó lớn nhất có thể.

\vspace{6pt}
\hangindent=13pt\hangafter=1
\textbf{8.} Tìm kích thước của một hình chữ nhật có diện tích $1000\text{ m}^2$ sao cho chu vi của nó nhỏ nhất có thể.

\vspace{6pt}
\hangindent=13pt\hangafter=1
\textbf{9.} Một mô hình được sử dụng cho năng suất $Y$ của một vụ mùa nông nghiệp như một hàm của mức nitơ $N$ trong đất (được đo bằng

\end{minipage}%
\hfill
% -------------------------- CỘT PHẢI --------------------------
\begin{minipage}[t]{0.48\textwidth}
\fontsize{9}{12.2}\selectfont

các đơn vị thích hợp) là
\[
Y = \frac{kN}{1 + N^2}
\]
trong đó $k$ là một hằng số dương. Mức nitơ nào mang lại năng suất tốt nhất?

\vspace{5pt}
\hangindent=18pt\hangafter=1
\textbf{10.} Tốc độ (tính bằng $\text{mg carbon}/\text{m}^3/\text{h}$) mà quá trình quang hợp diễn ra đối với một loài thực vật phù du được mô hình hóa bởi hàm số
\[
P = \frac{100I}{I^2 + I + 4}
\]
trong đó $I$ là cường độ ánh sáng (đo bằng hàng nghìn foot-candle). Cường độ ánh sáng bằng bao nhiêu thì $P$ đạt cực đại?

\vspace{5pt}
\hangindent=18pt\hangafter=1
\textbf{11.} Xét bài toán sau: một người nông dân có $750\text{ ft}$ hàng rào muốn rào một khu đất hình chữ nhật rồi chia nó thành bốn chuồng bằng các hàng rào song song với một cạnh của hình chữ nhật. Tổng diện tích lớn nhất có thể của bốn chuồng là bao nhiêu?
\begin{enumerate}[label=\textbf{(\alph*)}, leftmargin=18pt, labelsep=4pt, itemsep=2pt, parsep=0pt, topsep=2.5pt]
    \item Vẽ một số hình minh họa tình huống, một số hình có các chuồng nông và rộng, một số hình có các chuồng sâu và hẹp. Hãy tìm tổng diện tích của các cấu hình này. Liệu có vẻ như tồn tại một diện tích lớn nhất không? Nếu có, hãy ước lượng giá trị đó.
    \item Vẽ một hình minh họa tình huống tổng quát. Đưa vào các ký hiệu và gán nhãn cho hình vẽ bằng các ký hiệu của bạn.
    \item Viết một biểu thức cho diện tích tổng cộng.
    \item Sử dụng thông tin đã cho để viết một phương trình liên hệ giữa các biến.
    \item Sử dụng phần (d) để viết diện tích tổng cộng như một hàm của một biến duy nhất.
    \item Hoàn thành việc giải bài toán và so sánh câu trả lời với ước lượng của bạn ở phần (a).
\end{enumerate}

\vspace{5pt}
\hangindent=18pt\hangafter=1
\textbf{12.} Xét bài toán sau: một chiếc hộp không nắp được làm từ một tấm bìa các-tông hình vuông, cạnh rộng $3\text{ ft}$, bằng cách

\end{minipage}

\end{document}